% ============================================================
%  slide_ch2.tex — Chương II: Cơ sở lý thuyết
%  8 frames
% ============================================================
\section{Chương II: Cơ sở lý thuyết}

% ------ Frame 1: MDP Định nghĩa ------
\begin{frame}{Mô hình Quyết định Markov (MDP)}
  \begin{block}{Định nghĩa: $\mathcal{M} = (S,\; A,\; P,\; R,\; \gamma)$}
  \end{block}
  \vspace{4pt}
  \begin{columns}[T]
    \column{0.52\textwidth}
    \begin{tabular}{@{}lp{5.5cm}}
      $S$ & \textbf{Không gian trạng thái} --- tập trạng thái môi trường \\[4pt]
      $A$ & \textbf{Không gian hành động} --- tập quyết định của agent \\[4pt]
      $P(s'|s,a)$ & \textbf{Hàm chuyển trạng thái} --- xác suất chuyển \\[4pt]
      $R(s,a)$ & \textbf{Hàm phần thưởng} --- tín hiệu phản hồi \\[4pt]
      $\gamma\in(0,1]$ & \textbf{Hệ số chiết khấu} --- tầm nhìn dài hạn \\
    \end{tabular}

    \column{0.44\textwidth}
    \begin{alertblock}{Tính chất Markov}
      \[P(s_{t+1}\mid s_t, a_t) = P(s_{t+1}\mid s_0,\ldots,s_t,a_t)\]
      \centering Tương lai \textbf{chỉ phụ thuộc}\\vào trạng thái \textit{hiện tại}
    \end{alertblock}
    \vspace{6pt}
    \begin{block}{Chính sách tối ưu}
      \[\pi^* = \arg\max_\pi\; \mathbb{E}_\pi\!\left[\sum_{t=0}^{\infty}\gamma^t R_t\right]\]
    \end{block}
  \end{columns}
\end{frame}

% ------ Frame 2: Tại sao đèn giao thông là MDP ------
\begin{frame}{Tại sao điều khiển đèn giao thông là một MDP?}
  \begin{columns}[T]
    \column{0.52\textwidth}
    \begin{enumerate}
      \setlength{\itemsep}{6pt}
      \item \textbf{Quyết định tuần tự}\\
            Pha đèn tại $t$ ảnh hưởng trực tiếp trạng thái $t+1$
      \item \textbf{Tính ngẫu nhiên}\\
            Xe đến ngẫu nhiên, không thể dự đoán chính xác $s_{t+1}$
      \item \textbf{Tối ưu dài hạn}\\
            Cần giảm ùn tắc bền vững trong cả 1 giờ ($\gamma=0.99$)
      \item \textbf{Không biết mô hình vật lý}\\
            $P(s'|s,a)$ không tường minh → phải học từ dữ liệu
    \end{enumerate}

    \column{0.44\textwidth}
    \begin{block}{Ánh xạ bài toán $\Leftrightarrow$ MDP}
      \begin{center}
      \begin{tabular}{ll}
        \toprule
        \textbf{MDP} & \textbf{Giao thông}\\
        \midrule
        State $s_t$   & Hàng đợi, tốc độ\\
        Action $a_t$  & Pha đèn (0--3)\\
        Reward $r_t$  & $-$Mức ùn tắc\\
        Policy $\pi$  & Chiến lược điều khiển\\
        \bottomrule
      \end{tabular}
      \end{center}
    \end{block}
    \vspace{6pt}
    \begin{exampleblock}{}
      \centering$\gamma = 0.99$ → agent ưu tiên\\giảm ùn tắc dài hạn
    \end{exampleblock}
  \end{columns}
\end{frame}

% ------ Frame 3: Học tăng cường ------
\begin{frame}{Học tăng cường (Reinforcement Learning)}
  \begin{columns}[T]
    \column{0.52\textwidth}
    \textbf{Vòng lặp tương tác Agent -- Môi trường:}
    \vspace{6pt}

    \begin{tikzpicture}[scale=0.9, font=\small,
      box/.style={draw, rounded corners, minimum width=2cm, minimum height=0.7cm, fill=lightblue}]
      \node[box, fill=vnblue!25] (agent) at (0,2) {\textbf{Agent}};
      \node[box, fill=techgreen!25] (env)   at (4,2) {\textbf{Môi trường}};
      \draw[->, thick, vnblue]  (agent.east) to[bend left=25]
            node[above]{\small $a_t$ (pha đèn)} (env.west);
      \draw[->, thick, softorange] (env.west) to[bend left=25]
            node[below]{\small $s_{t+1},\; r_t$} (agent.east);
    \end{tikzpicture}

    \vspace{8pt}
    \textbf{Mục tiêu:} tối đa hóa return tích lũy
    \[G_t = \sum_{k=0}^{\infty}\gamma^k r_{t+k+1}\]

    \column{0.44\textwidth}
    \begin{block}{Đặc điểm nổi bật của RL}
      \begin{itemize}
        \item Không cần nhãn tường minh như Supervised Learning
        \item Phần thưởng có thể bị trễ (delayed reward)
        \item Tối ưu hiệu quả \textit{dài hạn}
        \item Phù hợp bài toán điều khiển trong môi trường động
      \end{itemize}
    \end{block}
    \vspace{4pt}
    \begin{alertblock}{Thách thức}
      Cân bằng \textbf{khai thác} (exploitation) và \textbf{khám phá} (exploration)
    \end{alertblock}
  \end{columns}
\end{frame}

% ------ Frame 4: Q-Learning & Bellman ------
\begin{frame}{Q-Learning và Phương trình Bellman}
  \begin{block}{Hàm giá trị hành động $Q^\pi(s,a)$}
    \[Q^\pi(s,a) = \mathbb{E}_\pi\bigl[G_t \mid s_t=s,\; a_t=a\bigr]
      = \mathbb{E}_\pi\!\left[\sum_{k=0}^{\infty}\gamma^k r_{t+k+1}\;\middle|\;s_t=s,\;a_t=a\right]\]
  \end{block}
  \vspace{6pt}
  \begin{columns}[T]
    \column{0.50\textwidth}
    \begin{block}{Phương trình Bellman tối ưu}
      \[Q^*(s,a) = \mathbb{E}\!\left[r + \gamma\max_{a'}Q^*(s',a')\;\middle|\;s,a\right]\]
      \[\Rightarrow\quad\pi^*(s) = \arg\max_a Q^*(s,a)\]
    \end{block}
    \vspace{4pt}
    \textbf{Q-Learning cập nhật:}
    \[Q(s,a)\;\leftarrow\; Q(s,a)+\alpha\Bigl[r+\gamma\max_{a'}Q(s',a')-Q(s,a)\Bigr]\]

    \column{0.46\textwidth}
    \begin{alertblock}{Giới hạn Q-table truyền thống}
      \begin{itemize}
        \item State space \textbf{liên tục} và \textbf{32 chiều}
        \item Bảng Q có kích thước vô hạn → \textbf{bất khả thi}
        \item[$\Rightarrow$] Cần xấp xỉ bằng \textbf{mạng nơ-ron sâu}
      \end{itemize}
    \end{alertblock}
    \vspace{6pt}
    \begin{exampleblock}{Giải pháp: Deep Q-Network}
      \[Q(s,a) \approx Q(s,a;\,\theta)\]
      $\theta$ = trọng số mạng nơ-ron
    \end{exampleblock}
  \end{columns}
\end{frame}

% ------ Frame 5: DQN ------
\begin{frame}{Deep Q-Network (DQN)}
  \begin{columns}[T]
    \column{0.52\textwidth}
    \textbf{Hàm mất mát (Huber / Smooth L1):}
    \[\mathcal{L}(\theta) = \mathbb{E}_{(s,a,r,s')\sim\mathcal{B}}\bigl[\mathrm{SmoothL1}(Q(s,a;\theta)-y)\bigr]\]

    \textbf{TD Target:}
    \[y = r + \gamma\max_{a'}Q(s',a';\theta^-)\]

    \vspace{6pt}
    \begin{block}{Smooth L1 Loss (Huber)}
      \[\mathrm{SmoothL1}(x)=\begin{cases}\tfrac{1}{2}x^2 & |x|\le 1\\ |x|-\tfrac{1}{2} & |x|>1\end{cases}\]
      Ít nhạy cảm với \textit{outlier} hơn MSE
    \end{block}

    \column{0.44\textwidth}
    \begin{block}{Kỹ thuật 1: Experience Replay}
      \begin{itemize}
        \item Lưu $(s,a,r,s',d)$ vào replay buffer
        \item Sample minibatch \textbf{ngẫu nhiên}
        \item Phá vỡ tương quan dữ liệu liên tiếp
        \item Tái sử dụng kinh nghiệm nhiều lần
      \end{itemize}
    \end{block}
    \vspace{6pt}
    \begin{block}{Kỹ thuật 2: Target Network $\theta^-$}
      \begin{itemize}
        \item Mạng $\theta^-$ tách biệt với mạng online $\theta$
        \item Cập nhật mỗi $N$ bước (hard copy)
        \item Tránh \textit{moving target problem}
        \item Ổn định quá trình hội tụ
      \end{itemize}
    \end{block}
  \end{columns}
\end{frame}

% ------ Frame 6: ε-Greedy ------
\begin{frame}{Chiến lược khám phá $\varepsilon$-Greedy}
  \begin{columns}[T]
    \column{0.55\textwidth}
    \[a_t = \begin{cases}
      \text{hành động ngẫu nhiên} & \text{xác suất } \varepsilon \\
      \arg\max_a Q(s_t,a;\theta) & \text{xác suất } 1-\varepsilon
    \end{cases}\]

    \vspace{8pt}
    \textbf{Lịch giảm tuyến tính:}
    \[\varepsilon(t) = \varepsilon_{\text{start}} + \frac{\min(t,T)}{T}\cdot(\varepsilon_{\text{end}}-\varepsilon_{\text{start}})\]

    \begin{center}
    \begin{tabular}{lc}
      \toprule
      \textbf{Tham số} & \textbf{Giá trị}\\
      \midrule
      $\varepsilon_{\text{start}}$ & 1.0 (toàn khám phá)\\
      $\varepsilon_{\text{end}}$   & 0.05 (5\% khám phá)\\
      $T$ (số bước giảm)           & 100.000\\
      \bottomrule
    \end{tabular}
    \end{center}

    \column{0.42\textwidth}
    \begin{tikzpicture}[scale=0.85]
      \draw[->, thick] (0,0) -- (5,0) node[right]{\small Bước $t$};
      \draw[->, thick] (0,0) -- (0,3.2) node[above]{\small $\varepsilon$};
      % epsilon line
      \draw[thick, vnblue, line width=2pt] (0,2.8) -- (2.8,0.4) -- (4.5,0.4);
      % labels
      \node[left]{\small 1.0} at (0,2.8);
      \node[left] at (0,0.4) {\small 0.05};
      \node[below] at (2.8,0) {\small 100K};
      \node[below] at (4.5,0) {\small 200K};
      % dashed
      \draw[dashed, gray] (2.8,0) -- (2.8,0.4);
      \draw[dashed, gray] (0,0.4) -- (2.8,0.4);
      % annotations
      \node[align=center, fill=vnblue!15, rounded corners, font=\tiny] at (1.4,1.8)
        {Khám phá nhiều\\(thu thập kinh nghiệm)};
      \node[align=center, fill=techgreen!15, rounded corners, font=\tiny] at (3.7,1.2)
        {Khai thác\\(tinh chỉnh)};
    \end{tikzpicture}
  \end{columns}
\end{frame}

% ------ Frame 7: Double DQN ------
\begin{frame}{Double DQN --- Khắc phục Overestimation Bias}
  \begin{columns}[T]
    \column{0.48\textwidth}
    \begin{alertblock}{Vấn đề DQN tiêu chuẩn}
      Cùng 1 mạng: vừa \textbf{chọn} vừa \textbf{đánh giá} hành động\\[4pt]
      \[y_{\text{DQN}} = r + \gamma\underbrace{\max_{a'}Q(s',a';\theta^-)}_{\text{chọn + đánh giá = 1 mạng}}\]
      $\Rightarrow$ Ước lượng Q quá cao (\textit{overestimation bias})
    \end{alertblock}

    \column{0.48\textwidth}
    \begin{exampleblock}{Giải pháp Double DQN}
      \textbf{Tách biệt hai vai trò:}
      \begin{align*}
        a^* &= \arg\max_{a'}Q(s',a';\alert{\theta})\quad\text{(online chọn)}\\
        y   &= r + \gamma\, Q(s',a^*;\alert{\theta^-})\quad\text{(target đánh giá)}
      \end{align*}
      $\Rightarrow$ Giảm overestimation, Q chính xác hơn
    \end{exampleblock}
  \end{columns}

  \vspace{8pt}
  \begin{center}
  \begin{tabular}{lcc}
    \toprule
    & \textbf{DQN tiêu chuẩn} & \textbf{Double DQN}\\
    \midrule
    Chọn $a^*$    & mạng target $\theta^-$ & mạng online $\theta$\\
    Đánh giá $Q$  & mạng target $\theta^-$ & mạng target $\theta^-$\\
    Overestimation & Cao & \textcolor{techgreen}{\textbf{Giảm đáng kể}}\\
    \bottomrule
  \end{tabular}
  \end{center}
\end{frame}

% ------ Frame 8: Dueling Network ------
\begin{frame}{Dueling Network Architecture}
  \begin{columns}[T]
    \column{0.50\textwidth}
    \textbf{Tách biệt 2 luồng:}
    \begin{itemize}
      \item \textbf{Value stream} $V(s;\theta_v)$:\\
            Giá trị của trạng thái $s$ --- bất kể hành động nào
      \item \textbf{Advantage stream} $A(s,a;\theta_a)$:\\
            Lợi thế tương đối của hành động $a$ tại $s$
    \end{itemize}
    \vspace{6pt}
    \textbf{Kết hợp (đảm bảo identifiability):}
    \[Q(s,a;\theta) = V(s) + \!\left[A(s,a) - \frac{1}{|A|}\sum_{a'}A(s,a')\right]\]
    \vspace{4pt}
    \begin{exampleblock}{Ưu điểm cho giao thông}
      Nhiều tình huống: giá trị \textit{trạng thái} quan trọng hơn phân biệt hành động → $V(s)$ học hiệu quả hơn
    \end{exampleblock}

    \column{0.46\textwidth}
    \begin{center}
    \begin{tikzpicture}[scale=0.82, font=\small,
      lyr/.style={draw, rounded corners, minimum width=3cm, minimum height=0.55cm, align=center}]

      \node[lyr, fill=vnblue!20]     (inp) at (0,4.8) {Input $s$ \;(32d)};
      \node[lyr, fill=vnblue!30]     (b1)  at (0,3.8) {FC + LayerNorm + ReLU \;(128)};
      \node[lyr, fill=vnblue!30]     (b2)  at (0,2.8) {FC + LayerNorm + ReLU \;(128)};
      \node[lyr, fill=techgreen!35]  (val) at (-1.6,1.6) {Value $V$ \;(1)};
      \node[lyr, fill=softorange!40] (adv) at (1.6,1.6)  {Advantage $A$ \;(4)};
      \node[lyr, fill=gray!25]       (q)   at (0,0.5) {$Q(s,a)$ \;(4 actions)};

      \draw[->, vnblue]        (inp) -- (b1);
      \draw[->, vnblue]        (b1)  -- (b2);
      \draw[->, techgreen]     (b2)  -- (val);
      \draw[->, softorange]    (b2)  -- (adv);
      \draw[->, techgreen]     (val) -- (q);
      \draw[->, softorange]    (adv) -- (q);
    \end{tikzpicture}
    \end{center}
  \end{columns}
\end{frame}
